\section{algorithm}
\subsection{L\'{e}vy-score-based Monte Carlo}

By virtue of Theorems \ref{theo:well-posedness}, \ref{theo:p_infty} and \ref{thm:gap_from_shell_cheeger}, the jump-diffusion SDE \eqref{eqn:sdeLevy2} provides a theoretically rigorous framework for sampling the Boltzmann distribution associated with \eqref{eqn:sde}. As a preliminary step, an appropriate L\'{e}vy process and its corresponding L\'{e}vy measure $\nu$ must be specified. In practice, this jump process must be amenable to efficient numerical implementation. One of the most computationally tractable choices is the compound Poisson process \cite{applebaum2009levy}:
\begin{equation}
    L^\mathrm{cp}_t=\sum_{i=1}^{N_t}A_i\delta(t-t_i),
\end{equation}
where $N_t$ denotes a Poisson process with intensity parameter $\lambda_0$, and the jump sizes $A_i$ are independent and identically distributed (i.i.d.) random variables drawn from a specified distribution $\nu_J$. The parameter $\lambda_0$ governs the frequency of jump events, while $\nu_J$ dictates the distribution of the jump magnitudes. In this work, we adopt the compound Poisson process as the primary jump mechanism for the proposed sampling algorithm, see Algorithm \ref{algorithm2}.

\begin{algorithm}
\caption{L\'{e}vy-score-based Monte Carlo sampling Boltzmann distribution}
\label{algorithm2}
\begin{algorithmic}[1]
    \STATE \textbf{Input:} An initial state $Z_{t_0}=z_0$, time step size $\Delta t$, total simulation time $T$, Poisson intensity parameter $\lambda_0$, and an appropriate jump size distribution $\nu_J$.
    \FOR{$k=0, 1, \dots, T/\Delta t - 1$}
        \STATE Sample the standard normal variable $\xi_k\sim \mathcal{N}(0,1)$;
        \STATE Sample the number of jumps $N_{\Delta t_k}\sim \mathrm{Po}(\lambda_0 \Delta t)$, and sample the jump sizes $A_i\sim\nu_J$ for $i=1,\dots,N_{\Delta t_k}$;
        \STATE Propagate the system state: 
        \begin{align*}
            Z_{t_{k+1}}=& Z_{t_{k}}+ \Delta t \left[ -\nabla V(Z_{t_{k}}) -\int_0^1\int_{\mathbb{R}^d-\{0\}} r\exp\left\{-\frac{2V(Z_{t_{k}}- \theta r)}{\sigma^2} + \frac{2V(Z_{t_{k}})}{\sigma^2}\right\} \nu(\d r)\d\theta \right] \\
            & \quad + \sigma\sqrt{\Delta t}\xi_k + \sum_{i=1}^{N_{\Delta t_k}}A_i;
        \end{align*}
        \STATE Update the time variable $t_{k+1}=t_{k}+\Delta t$;
    \ENDFOR
    \STATE \textbf{Output:} The trajectory samples $\{Z_{t_k}\}_{k=1}^{T/\Delta t}$ and the empirical distribution $p_N(x):=\frac{1}{T}\sum_{k=0}^{T/\Delta t}\delta(x-Z_{t_k})$.
\end{algorithmic}
\end{algorithm}

\begin{remark}
    Determining the optimal choices for the Poisson intensity parameter $\lambda_0$ and the jump size distribution $\nu_J$ to maximize the efficiency of sampling the Boltzmann distribution $p_\infty$ in \eqref{p_infty} remains an open and pertinent question. The optimal configuration is inherently dependent on the specific landscape of the potential $V$. Because the energy barriers separating local minima profoundly limit the efficiency of standard Monte Carlo algorithms, $\nu_J$ should ideally be tailored such that the spatial jump magnitudes are sufficient to traverse the distances between distinct wells. If the typical jump sizes are too small to bypass these regional basins of attraction, the non-local process merely introduces localized fluctuations analogous to an auxiliary continuous diffusion, failing to facilitate the desired macroscopic state transitions. A rigorous theoretical and numerical investigation of these optimal parameter selections is reserved for future work.
\end{remark}


